\section{Related Work}
\label{sec:related}

\paragraph{Diffusion Language Models}
While diffusion models have achieved remarkable success in the visual domain \citep{songscore, ho2020denoising}, their application to language has been limited, partly due to text's discrete nature. Initial approaches attempted to learn continuous diffusion models over textual latents \citep{austin2021structured, gulrajani2023likelihood}, but faced challenges with scalability and discretization. Masked diffusion has been established as a specific instance of discrete diffusion \citep{austin2021structured, sahoo2024simple, shi2024simplified, ou2024your,nie2024scaling}, with recent efforts scaling these models significantly. DiffuLLaMA \citep{gong2025scaling} extended this approach by initializing masked diffusion language models with pretrained LLaMA weights. \citet{yediffusion} explored how diffusion language models can generate chain-of-thought reasoning, and complex reasoning tasks on smaller-scale models~\citep{ye2024beyond}, highlighting their advantages over autoregressive models in reversal tasks, though their traces lacked self-correction capabilities. \citet{arriola2025block} proposed Block Diffusion, a hybrid approach that models sequences block-by-block while applying diffusion within each block, allowing flexible length generation and improving inference efficiency with kv-caching. Recently, LLaDA \citep{nie2025largelanguagediffusionmodels} and Dream \citep{dream2025} demonstrated that large diffusion language models can achieve performance comparable to similarly-sized autoregressive alternatives, but have not yet been enhanced through reinforcement learning. To the best of our knowledge, we are the first to demonstrate the efficacy of policy gradient-based reinforcement learning algorithms on large diffusion language models. 

\paragraph{Improving Reasoning Abilities of LLMs through SFT and RL}
Approaches to enhance reasoning capabilities in large language models generally fall into two categories: supervised finetuning and reinforcement learning. SFT on high-quality reasoning traces~\citep{yu2023metamath, numina_math_datasets, pasteropenwebmath} has shown promising results, while fewer but carefully curated reasoning datasets~\citep{ye2025limoreasoning, muennighoff2025s1, zhou2023lima} can outperform larger datasets. \citet{chu2025sft} demonstrate that SFT-based reasoning often relies on memorization rather than generalization, while RL methods achieve better transfer to novel scenarios, particularly when intermediate reasoning steps are difficult to supervise. Recently, algorithms like GRPO~\citep{guo2025deepseek, shao2024deepseekmath} enable efficient training by estimating advantages from group scores without requiring additional critic models as in PPO. \citet{guo2025deepseek} demonstrate that strong reasoning capabilities can emerge through RL even without SFT (DeepSeek-R1-Zero), producing long reasoning traces with self-reflection and verification steps that significantly improve performance on mathematical tasks. The development of strong reasoning models like R1 has in turn sparked renewed interest in SFT for smaller models using distilled reasoning traces from these expert reasoners. Datasets like OpenThoughts~\citep{openthoughts} and OpenR1-Math\footnote{\url{https://huggingface.co/datasets/open-r1/OpenR1-Math-220k}}, which contain reasoning traces from DeepSeek R1, enable smaller models to learn step-by-step problem-solving from expert demonstrations. For RL in discrete diffusion models, prior work by \citet{zekri2025fine} proposed a policy gradient framework using concrete score matching, but it relies on gradient-flow computations and does not target masked objectives. In contrast, our method is tailored to masked dLLMs with efficient policy gradient calculation and improved learning efficiency through random masking. Our work is among the first to explore improving reasoning in diffusion-based LLMs via both SFT and RL.


